In chapter 3 we studied neural networks in the lazy regime.
In chapter 4 we studied how to get neural networks out of the lazy regime
and into the rich regime. 
However, we have yet to talk about what actually goes on in the rich,
feature learning regime. 
The problem is that there really is no consensus (as of 2026) among researchers
on how to think about feature learning. 
Different factions utilize different approaches and have varying beliefs about which open 
questions matter most.
In this chapter, we'll explore a handful of research programs---what kind of questions
they seek to answer, what successes they've had, their strengths and weaknesses, etc.

\section{Research programs under the umbrella of learning mechanics}

\subsection{Simplifying assumptions and toy models}

\section{Modeling structure in natural data}

\section{Optimization dynamics}

\section{Applications of statistical physics}

\section{Mechanistic Interpretability}

\section{A case for optimism}